\section{Limitations}
\textbf{No human baseline.} Recruitment of qualified crystallographers failed. We therefore make no claim
that the renders are solvable by humans; the oracle stands in for the information question only, and
``adequate'' means recoverable by a geometric reader.

\textbf{What the oracle does not license.} It reads ground-truth coordinates, not pixels. It bounds what is
recoverable given ideal extraction and says nothing about achievable extraction. The extraction share of
the ceiling-to-model gap is bounded from neither side: $R_1$ assumes perfect extraction and $R_2$'s
extractor is too weak to be informative. A detector at recall above 0.9 is the most valuable missing
measurement.

\textbf{Sample size.} All model numbers rest on two samples of 210; the minimum detectable paired
difference there is near 0.09, so adjacent leaderboard rows are not separable. The model arms were not
re-run at $n = 1933$ (that is $19{,}500$ calls on a 500-structure core), so the ceiling-to-model gap at
scale is bounded on the oracle side only. The oracle and both classical baselines are complete on the full
1933.

\textbf{Corrections and non-recoveries.} Three published values did not reproduce from their own records
and are documented as non-recoveries rather than silently replaced: an earlier cue-sufficiency classifier's
exact rule, an oracle view-count curve, and the tabular classifier's published count (twelve defensible
readings of its recorded protocol span 183 to 188 without reaching the published 186). A previously
published minimum-detectable-difference for the render-convention null was computed with an unpaired
approximation and is retracted in favour of the per-comparison analysis of Section~\ref{sec:conventions},
which makes that null weaker rather than stronger. The full retraction ledger is in the supplementary
(Appendix~\ref{app:corrections} indexes it).
